% Options for packages loaded elsewhere
\PassOptionsToPackage{unicode}{hyperref}
\PassOptionsToPackage{hyphens}{url}
\documentclass[
]{article}
\usepackage{xcolor}
\usepackage{amsmath,amssymb}
\setcounter{secnumdepth}{-\maxdimen} % remove section numbering
\usepackage{iftex}
\ifPDFTeX
  \usepackage[T1]{fontenc}
  \usepackage[utf8]{inputenc}
  \usepackage{textcomp} % provide euro and other symbols
\else % if luatex or xetex
  \usepackage{unicode-math} % this also loads fontspec
  \defaultfontfeatures{Scale=MatchLowercase}
  \defaultfontfeatures[\rmfamily]{Ligatures=TeX,Scale=1}
\fi
\usepackage{lmodern}
\ifPDFTeX\else
  % xetex/luatex font selection
\fi
% Use upquote if available, for straight quotes in verbatim environments
\IfFileExists{upquote.sty}{\usepackage{upquote}}{}
\IfFileExists{microtype.sty}{% use microtype if available
  \usepackage[]{microtype}
  \UseMicrotypeSet[protrusion]{basicmath} % disable protrusion for tt fonts
}{}
\makeatletter
\@ifundefined{KOMAClassName}{% if non-KOMA class
  \IfFileExists{parskip.sty}{%
    \usepackage{parskip}
  }{% else
    \setlength{\parindent}{0pt}
    \setlength{\parskip}{6pt plus 2pt minus 1pt}}
}{% if KOMA class
  \KOMAoptions{parskip=half}}
\makeatother
\usepackage{longtable,booktabs,array}
\newcounter{none} % for unnumbered tables
\usepackage{calc} % for calculating minipage widths
% Correct order of tables after \paragraph or \subparagraph
\usepackage{etoolbox}
\makeatletter
\patchcmd\longtable{\par}{\if@noskipsec\mbox{}\fi\par}{}{}
\makeatother
% Allow footnotes in longtable head/foot
\IfFileExists{footnotehyper.sty}{\usepackage{footnotehyper}}{\usepackage{footnote}}
\makesavenoteenv{longtable}
\setlength{\emergencystretch}{3em} % prevent overfull lines
\providecommand{\tightlist}{%
  \setlength{\itemsep}{0pt}\setlength{\parskip}{0pt}}
\usepackage{bookmark}
\IfFileExists{xurl.sty}{\usepackage{xurl}}{} % add URL line breaks if available
\urlstyle{same}
\hypersetup{
  hidelinks,
  pdfcreator={LaTeX via pandoc}}

\author{}
\date{}

\begin{document}

\section{Explicit tails for the original test at time
1/32}\label{explicit-tails-for-the-original-test-at-time-132}

Date: 2026-09-23. This is an independent derivation for the original
source and filter in ER2, ER27, ER34, and ER39--40 of
\texttt{../endpoint\_resonance\_20260923/ENDPOINT\_RESONANCE\_AND\_PRIME\_CLASS.tex}.

\subsection{Source and reading record}\label{source-and-reading-record}

The source file was read in full. The calculations below use its exact
theta density (ER2), actual test (ER27), differential filter and
convolution (ER34), and complete arithmetic and Gamma pairing
(ER39--40). The human source identified there is Brad Rodgers and
Terence Tao, \emph{The de Bruijn--Newman constant is non-negative},
\href{https://arxiv.org/abs/1801.05914v5}{arXiv:1801.05914v5},
original-source equation identifiers \texttt{phidef} and \texttt{htdef}.
This independent derivation read the programme file containing those
formulas; it does not claim a new reading of the original author source
archive. The explicit-formula convention is the one already proved in
ER39--40, with Alain Connes's
\href{https://arxiv.org/abs/math/9811068v1}{original paper}, Appendix
II, cited there.

No hypothesis on the positions of zeta zeros, no numerical zero list,
and no completed-function substitution enters any bound below. The
displayed inequalities majorize the original functions for error
control. They do not change the functions used in the pairing.

\subsection{1. Original definitions and uniform derivative
moments}\label{original-definitions-and-uniform-derivative-moments}

Throughout this note the heat time is exactly

{[} t=\frac1{32},\qquad \frac1{4t}=8,\qquad \frac1{2t}=16. {]}

Retain the entire original density

{[} \Phi(u)=\sum\_\{n=1\}\^{}\{\infty\}
\bigl(2\pi\textsuperscript{2n}4e\^{}\{9u\}-3\pi n\textsuperscript{2e}\{5u\}\bigr)
e\^{}\{-\pi n\textsuperscript{2e}\{4u\}\},\qquad u\ge0, {]}

and define

{[}
H(x)=\int\_0\^{}\infty e\textsuperscript{\{u}2/32\}\Phi(u)\cos(xu),du,
\qquad M\_j=\int\_0\^{}\infty u\^{}j
e\textsuperscript{\{u}2/32\}\Phi(u),du \quad(j\ge0). \tag{TB1} {]}

Each original summand is strictly positive because
(2\pi n\textsuperscript{2e}\{4u\}\textgreater3). In an upper bound it is
therefore legitimate to use

{[}
0\textless{}\Phi(u)\le2\pi\textsuperscript{2\sum\_\{n\ge1\}n}4e\textsuperscript{\{9u\}e}\{-\pi n\textsuperscript{2e}\{4u\}\}.
\tag{TB2} {]}

This inequality retains the negative summand in the definition TB1; it
records its sign explicitly when estimating the full sum. For every
(u\ge0), the exponential series gives

{[} e\textsuperscript{\{4u\}\ge1+4u+8u}2. {]}

Since (\pi\textgreater3), (\pi\^{}2\textless10), and
(8\pi n\^{}2-1/32\textgreater0),

{[}

\begin{aligned}
\frac{u^2}{32}+9u-\pi n^2e^{4u}
&\le-\pi n^2-(4\pi n^2-9)u-(8\pi n^2-1/32)u^2\\
&\le-3n^2-(12n^2-9)u.
\end{aligned}

{]}

Tonelli's theorem applies to the nonnegative majorant, and hence

{[}

\begin{aligned}
M_j
&\le20j!\sum_{n\ge1}\frac{n^4e^{-3n^2}}{(12n^2-9)^{j+1}}\\
&\le\frac{20j!}{3^{j+1}}\sum_{n\ge1}n^{2-2j}e^{-3n^2}
\le\frac{20j!}{3^{j+1}}\sum_{n\ge1}n^2e^{-3n^2}.
\end{aligned}
\tag{TB3}

{]}

Here (12n\textsuperscript{2-9\ge3n}2). The elementary estimate

{[}
\sum\_\{n\ge1\}n\textsuperscript{2e}\{-3n\^{}2\}\textless{}\frac1{20}
\tag{TB4} {]}

can be checked entirely with rational arithmetic. Indeed,

{[} e\textsuperscript{3\textgreater{}\sum\_\{j=0\}}\{10\}\frac{3^j}{j!}
=\frac{899569}{44800}\textgreater{}\frac{2007}{100}\textgreater20. {]}

For (n\ge2), (n\^{}2\ge1+3(n-1)). Put (q=1/8000). Then

{[}

\begin{aligned}
\sum_{n\ge1}n^2e^{-3n^2}
&<\frac{100}{2007}+\frac1{20}\sum_{n\ge2}n^2q^{n-1}\\
&=\frac{100}{2007}+\frac1{20}
\left(\frac{1+q}{(1-q)^3}-1\right)\\
&=\frac{1024129791832007}{20543974083319860}<\frac1{20}.
\end{aligned}

{]}

The rational generating-function identity follows by differentiating
(\sum\_\{n\ge0\}q\textsuperscript{n=(1-q)}\{-1\}) twice within
(\textbar q\textbar\textless1). Consequently

{[} \boxed{M_j\le\frac{j!}{3^{j+1}},\qquad
|H^{(j)}(x)|\le M_j\quad(x\in\mathbb R, j\ge0).} \tag{TB5} {]}

The same integrable bounds justify differentiating TB1 under the
integral sign. In particular, (H) is real and even,
(M\_0=H(0)\textgreater0), and (H'\,'(0)=-M\_2\textless0).

\subsection{2. The actual filtered function and its
derivatives}\label{the-actual-filtered-function-and-its-derivatives}

Retain the functions

{[} h(x)=e\textsuperscript{\{-8x}2\}H(x)\cos(16x),\qquad
f(x)=h'\,'(x)-\frac14h(x),\qquad k(x)=\int\_\{\mathbb R\}f(y)f(x-y),dy.
\tag{TB6} {]}

Direct differentiation, with the (-h/4) term included, gives the exact
formula

{[}

\begin{aligned}
f(x)=e^{-8x^2}\bigl\{&
[(256x^2-1089/4)H(x)-32xH'(x)+H''(x)]\cos(16x)\\
&+[512xH(x)-32H'(x)]\sin(16x)\bigr\}.
\end{aligned}
\tag{TB7}

{]}

Thus, with (r=\textbar x\textbar), TB5 implies

{[}

\begin{aligned}
|f(x)|
&\le e^{-8x^2}
\left[(256r^2+512r+1089/4)M_0+(32r+32)M_1+M_2\right]\\
&\le e^{-8x^2}\left(\frac{10193}{108}+\frac{1568}{9}r+\frac{256}{3}r^2\right)\\
&\le(95+175r+86r^2)e^{-8x^2}
\le(183+174x^2)e^{-8x^2}.
\end{aligned}
\tag{TB8}

{]}

The last inequality uses (175r\le(175/2)(1+r\^{}2)). Also,
(re\textsuperscript{\{-4r}2\}\le1/4) and
(r\textsuperscript{2e}\{-4r\^{}2\}\le1/10), because (e\textgreater5/2).
Therefore

{[} \boxed{|f(x)|\le148e^{-4x^2}.} \tag{TB9} {]}

Here is a finite formula giving all required derivative bounds without
omitting any derivative of (H). Define positive polynomials

{[} T\_m(r)=\sum\_\{a=0\}\^{}\{\lfloor m/2\rfloor\}
\frac{m!\,8^a\,16^{m-2a}}{a!(m-2a)!}(1+r)\^{}\{m-2a\}, {]}

{[} S\_m(r)=\sum\emph{\{\ell=0\}\^{}\{m\}\binom m\ell
\frac{\ell!}{3^{\ell+1}}T}\{m-\ell\}(r),\qquad
B\_j(r)=S\_\{j+2\}(r)+\frac14S\_j(r). \tag{TB10} {]}

To prove their bounds, let (a(x)=e\textsuperscript{\{-8x}2+16ix\}). Its
exact derivative polynomial is

{[} a\^{}\{(m)\}(x)=a(x) \sum\_\{b=0\}\^{}\{\lfloor m/2\rfloor\}
\frac{m!(-8)^b}{b!(m-2b)!}(-16x+16i)\^{}\{m-2b\}. \tag{TB11} {]}

This follows by taking the coefficient of (z\^{}m) in
(a(x+z)=a(x)\exp[(-16x+16i)z-8z^2]). Since
(\textbar-16x+16i\textbar{}\le16(1+r)), TB11 gives
(\textbar a\^{}\{(m)\}(x)\textbar{}\le e\textsuperscript{\{-8x}2\}T\_m(r)).
Leibniz's rule, TB5, and (h=\Re(aH)) prove

{[}
\textbar h\^{}\{(m)\}(x)\textbar{}\le e\textsuperscript{\{-8x}2\}S\_m(r),\qquad
\boxed{|f^{(j)}(x)|\le e^{-8x^2}B_j(r).} \tag{TB12} {]}

For convenient integer ceilings,

{[} \boxed{|f^{(j)}(x)|\le C_j e^{-4x^2}\quad(0\le j\le4),} \tag{TB13}
{]}

with the following constants:

{\def\LTcaptype{none} % do not increment counter
\begin{longtable}[]{@{}rr@{}}
\toprule\noalign{}
(j) & (C\_j) \\
\midrule\noalign{}
\endhead
\bottomrule\noalign{}
\endlastfoot
0 & 148 \\
1 & 3,350 \\
2 & 81,000 \\
3 & 2,050,000 \\
4 & 54,500,000 \\
\end{longtable}
}

For a reproducible rational verification, if
(B\_j(r)=\sum\emph{pb}\{jp\}r\^{}p), use

{[} (d\_0,d\_1,d\_2,d\_3,d\_4,d\_5,d\_6)
=\left(1,\frac14,\frac1{10},\frac3{50},\frac1{25},\frac1{32},\frac{27}{1000}\right).
{]}

For (0\le p\le6), (r\textsuperscript{pe}\{-4r\^{}2\}\le d\_p): the
maximum for (p\textgreater0) is ((p/(8e))\^{}\{p/2\}), and the displayed
rational ceilings follow from (e\textgreater5/2), including
(\sqrt{3/20}\textless2/5) for (p=3). Substituting TB10 gives the exact
upper sums

{[} \left(\sum\emph{pb}\{jp\}d\_p\right)\_\{j=0\}\^{}4 =\left(
\frac{79093}{540}, \frac{9042119}{2700}, \frac{327148231}{4050},
\frac{24888564437}{12150}, \frac{1653729054454}{30375} \right). {]}

Each is less than its corresponding (C\_j), proving the table.

\subsection{3. Convolution bounds with the original Gaussian exponent
retained}\label{convolution-bounds-with-the-original-gaussian-exponent-retained}

The identity

{[} y\textsuperscript{2+(x-y)}2=\frac{x^2}{2}+2(y-x/2)\^{}2 \tag{TB14}
{]}

and TB9 first give

{[}
\textbar k(x)\textbar{}\le148\textsuperscript{2\sqrt{\frac\pi8},e}\{-2x\^{}2\}
\textless14000e\textsuperscript{\{-2x}2\}. \tag{TB15} {]}

A stronger tail bound follows by convolving the final polynomial in TB8.
Put (a=183), (b=174), and (z=y-x/2). Integrating all terms, including
the negative quadratic contribution of (y\textsuperscript{2(x-y)}2),
gives

{[}

\begin{aligned}
|k(x)|
&\le\frac{\sqrt\pi}{4}e^{-4x^2}
\left[a^2+ab\left(\frac{x^2}{2}+\frac1{16}\right)
+b^2\left(\frac{x^4}{16}-\frac{x^2}{64}+\frac3{1024}\right)\right]\\
&=\frac{\sqrt\pi}{4}e^{-4x^2}
\left(\frac{9105363}{256}+\frac{247167}{16}x^2+\frac{7569}{4}x^4\right)\\
&\le Q(x)e^{-4x^2},\qquad
\boxed{Q(x)=16000+6900x^2+850x^4.}
\end{aligned}
\tag{TB16}

{]}

Indeed, for the Gaussian (e\textsuperscript{\{-16z}2\}), its total mass
is (\sqrt\pi/4), its second moment divided by that mass is (1/32), and
its fourth moment divided by that mass is (3/1024). This proves the
first line. The classical elementary bound
(\pi\textless22/7\textless256/81) implies (\sqrt\pi/4\textless4/9).
Multiplying the three coefficients by (4/9) gives respectively
(15807.921875), (6865.75), and (841), each below the corresponding
integer in (Q). The terminating decimals here are exact rational values.

Differentiation under the convolution integral is justified by TB13. One
may put any selected derivatives on either factor. In particular,

{[} k'=f'\emph{f,\quad k'\,`=f'}f',\quad
k'\,'\,`=f'`\emph{f',\quad k\^{}\{(4)\}=f'\,'}f'\,'. {]}

Using (\sqrt{\pi/8}\textless63/100) and TB13 gives

{[} \textbar k\^{}\{(j)\}(x)\textbar{}\le D\_j
e\textsuperscript{\{-2x}2\},\quad 0\le j\le4, \tag{TB17} {]}

where

{[} (D\_0,D\_1,D\_2,D\_3,D\_4)
=(14000,313000,7100000,171000000,4140000000). {]}

For example the pre-ceiling constants are ((63/100)148\^{}2),
((63/100)148\cdot3350), ((63/100)3350\^{}2), ((63/100)81000\cdot3350),
and ((63/100)81000\^{}2), respectively.

Both (f) and (k) are real and even. The original value satisfies

{[} k(0)=\int\_\{\mathbb R\}f(y)\^{}2,dy\textgreater0, {]}

because TB7 gives (f(0)=-M\_2-(1089/4)M\_0\textless0). Also,

{[} k(0)-k(x)=\frac12\int\_\{\mathbb R\}{[}f(y)-f(y-x){]}\^{}2,dy\ge0.
\tag{TB18} {]}

Expanding the square proves the identity, since translation preserves
the two square integrals and evenness converts the cross integral to
(k(x)). By the fundamental theorem of calculus, Cauchy--Schwarz, and
Tonelli,

{[}
0\le k(0)-k(x)\le\frac{x^2}{2}\int\_\{\mathbb R\}\textbar f'(y)\textbar\^{}2,dy
\le\frac{D_2x^2}{2}. \tag{TB19} {]}

The final inequality follows either directly from TB13 or from the
constants used in TB17.

\subsection{4. Prime tail, including every prime
power}\label{prime-tail-including-every-prime-power}

For an integer (N\ge3), define the actual omitted arithmetic term

{[}
P\_\{\textgreater N\}=2\sum\_\{n\textgreater N\}\frac{\Lambda(n)}{\sqrt n}k(\log n),
\quad \Lambda(p\^{}a)=\log p~(a\ge1),
\quad \Lambda(n)=0\text{ otherwise}. \tag{TB20} {]}

The elementary inequality (0\le\Lambda(n)\le\log n) holds at every
integer, including all prime powers. Put

{[} L=\log N,\quad c=8L-\frac12, \quad
\mathcal T\_p(L,c)=\sum\_\{j=0\}\^{}\{p\}\binom pj
\frac{L^{p-j}j!}{c^{j+1}}. \tag{TB21} {]}

Then

{[} \boxed{
|P_{>N}|\le
2e^{-4L^2+L/2}
\left[16000\mathcal T_1(L,c)
+6900\mathcal T_3(L,c)
+850\mathcal T_5(L,c)\right].} \tag{TB22} {]}

Here is the complete sum-to-integral justification. For (u\ge1), the
logarithmic derivative of (uQ(u)e\textsuperscript{\{-4u}2-u/2\}) is

{[} \frac1u+\frac{Q'(u)}{Q(u)}-8u-\frac12
\le\frac5u-8u-\frac12\textless0. {]}

The inequality (uQ'(u)\le4Q(u)) follows directly from the three positive
terms in TB16. Since (N\ge3\textgreater e), the nonnegative function
((\log x)x\^{}\{-1/2\}Q(\log x)e\^{}\{-4(\log x)\^{}2\}) decreases for
(x\ge N). Consequently TB16 and the integral comparison for a decreasing
function yield

{[} \textbar P\_\{\textgreater N\}\textbar{}
\le2\int\_N\^{}\infty(\log x)x\^{}\{-1/2\}Q(\log x)e\^{}\{-4(\log x)\^{}2\},dx
=2\int\_L\^{}\infty uQ(u)e\textsuperscript{\{-4u}2+u/2\},du. {]}

For (u=L+v), the exponent is exactly

{[} -4u\textsuperscript{2+u/2=-4L}2+L/2-cv-4v\^{}2. {]}

Bounding (e\textsuperscript{\{-4v}2\}\le1), expanding each of
((L+v)\textsuperscript{1,(L+v)}3,(L+v)\^{}5), and using
(\int\_0\^{}\infty v\textsuperscript{je}\{-cv\},dv=j!/c\^{}\{j+1\})
proves TB22. It also proves absolute convergence independently of any
information about zeros or prime density.

A useful completely rational cutoff consequence is

{[} \boxed{|P_{>64}|<10^{-21}.} \tag{TB23} {]}

To verify it, (e\textless11/4) implies
(e\textsuperscript{4\textless(11/4)}4\textless64), hence
(\log64\textgreater4). One elementary proof of the bound on (e) is

{[} e=\sum\emph{\{j=0\}\^{}3\frac1{j!}+\sum}\{j\ge4\}\frac1{j!}
\le\frac83+\frac1{24}\sum\_\{j\ge0\}5\^{}\{-j\}
=\frac{87}{32}\textless{}\frac{11}{4}. {]}

Increase the integral majorant's domain from ({[}\log64,\infty)) to
({[}4,\infty)). Applying TB22's polynomial integration with (L=4),
(c=63/2), its coefficient in front of (e\^{}\{-62\}) is exactly

{[}
2{[}16000\mathcal T\_1(4,63/2)+6900\mathcal T\_3(4,63/2)+850\mathcal T\_5(4,63/2){]}
=\frac{1882788609056000}{20841167403}\textless100000. {]}

Finally
(e\textsuperscript{\{62\}\textgreater e}\{60\}\textgreater20\textsuperscript{\{20\}\textgreater10}\{26\}),
proving TB23. Every original arithmetic term with (n\le64), including
powers of primes, remains in the finite calculation.

\subsection{5. Gamma integral: the combined numerator, endpoint, and
tail}\label{gamma-integral-the-combined-numerator-endpoint-and-tail}

For the original even convolution, ER40's Gamma integral is

{[} \int\_0\^{}\infty \frac{e^{-x}k(0)-e^{-x/4}k(x/2)}{1-e^{-x}},dx.
\tag{TB24} {]}

The prefactor in the full archimedean term remains
(-(\gamma\_\{\rm E\}+\log\pi)k(0)), exactly as in ER40. The combined
numerator can be written exactly as

{[}
(e\textsuperscript{\{-x\}-e}\{-x/4\})k(0)+e\^{}\{-x/4\}{[}k(0)-k(x/2){]}.
\tag{TB25} {]}

By TB19 the second bracket is between zero and (D\_2x\^{}2/8). Since
((e\textsuperscript{\{-x\}-e}\{-x/4\})/x\to-3/4) and
((1-e\^{}\{-x\})/x\to1), the integrand has the definite endpoint value

{[} \boxed{\lim_{x\downarrow0}
\frac{e^{-x}k(0)-e^{-x/4}k(x/2)}{1-e^{-x}}
=-\frac34k(0).} \tag{TB26} {]}

In particular no individual divergent integral has been used at zero.
More explicitly, for (0\textless x\le1),
(\textbar e\textsuperscript{\{-x\}-e}\{-x/4\}\textbar{}\le3x/4) and
(1-e\^{}\{-x\}\ge x/2), so the absolute value of the original integrand
is at most (3k(0)/2+D\_2x/4). This proves its local integrability with a
stated majorant.

For any (T\textgreater0), the actual tail admits the exact decomposition

{[}

\begin{aligned}
G_{>T}
&=\int_T^\infty
\frac{e^{-x}k(0)-e^{-x/4}k(x/2)}{1-e^{-x}}\,dx\\
&=-k(0)\log(1-e^{-T})-\mathcal R_T,\\
\mathcal R_T
&=\int_T^\infty\frac{e^{-x/4}k(x/2)}{1-e^{-x}}\,dx.
\end{aligned}
\tag{TB27}

{]}

The first term follows by the substitution (v=e\^{}\{-x\}). It can be
retained exactly in a finite-interval evaluation. Neither the sign of
(k(x/2)) nor that of (\mathcal R\_T) is assumed.

Using TB16, put (d=2T+1/4) and use (\mathcal T\_p) from TB21. Then

{[} \boxed{
|\mathcal R_T|\le
\frac{e^{-T^2-T/4}}{1-e^{-T}}
\left[16000\mathcal T_0(T,d)
+1725\mathcal T_2(T,d)
+\frac{425}{8}\mathcal T_4(T,d)\right].} \tag{TB28} {]}

Indeed,

{[} \textbar k(x/2)\textbar{}\le
\left(16000+1725x\textsuperscript{2+\frac{425}{8}x}4\right)e\textsuperscript{\{-x}2\},
{]}

and
((1-e\textsuperscript{\{-x\})}\{-1\}\le(1-e\textsuperscript{\{-T\})}\{-1\})
for (x\ge T). Under (x=T+v), the complete exponent is
(-T\textsuperscript{2-T/4-dv-v}2). Bounding its last factor by one and
integrating the expanded polynomial gives TB28.

In particular,

{[} \boxed{G_{>8}=-k(0)\log(1-e^{-8})-\mathcal R_8,
\qquad |\mathcal R_8|<10^{-23}.} \tag{TB29} {]}

For the rational check, (d=65/4), and the bracket in TB28 is

{[} 16000\mathcal T\_0(8,65/4)+1725\mathcal T\_2(8,65/4)
+\frac{425}{8}\mathcal T\_4(8,65/4) =\frac{1006957513344}{46411625}. {]}

Its double is less than (100000); also
((1-e\textsuperscript{\{-8\})}\{-1\}\textless2). Since
(e\textsuperscript{\{66\}\textgreater20}\{22\}\textgreater10\^{}\{28\}),
TB28 is less than (100000/10\textsuperscript{\{28\}=10}\{-23\}). Thus
the finite Gamma integral over ({[}0,8{]}), its endpoint value TB26, the
exact term (-k(0)\log(1-e\^{}\{-8\})), and the bound on (\mathcal R\_8)
keep all original archimedean contributions.

For a shorter but weaker bound, TB15 gives

{[} \textbar{}\mathcal R\_T\textbar{}\le
\frac{14000e^{-T^2/2-T/4}}{(1-e^{-T})(T+1/4)}. \tag{TB30} {]}

This follows from (x\^{}2/2+x/4\ge T\^{}2/2+T/4+(T+1/4)(x-T)). If the
exact first term in TB27 is not evaluated, add
(k(0){[}-\log(1-e\^{}\{-T\}){]}) to either TB28 or TB30 to bound
(\textbar G\_\{\textgreater T\}\textbar).

\subsection{6. Exact constant-plus-remainder comparison requested by the
receiving
calculation}\label{exact-constant-plus-remainder-comparison-requested-by-the-receiving-calculation}

This section verifies the root calculation's comparison with its full
residual. Put

{[} C=H(0)=M\_0,\qquad R(x)=H(x)-C, {]}

{[} f\_0(x)=Ce\textsuperscript{\{-8x}2\}
{[}(256x\^{}2-1089/4)\cos(16x)+512x\sin(16x){]},
\qquad e\_f(x)=f(x)-f\_0(x). \tag{TB31} {]}

The identity (H'(0)=0), Taylor's integral formula, and TB1 give

{[} \textbar R(x)\textbar{}\le\frac{M_2x^2}{2},\quad
\textbar R'(x)\textbar{}\le M\_2\textbar x\textbar,\quad
\textbar R'\,`(x)\textbar{}\le M\_2,\quad
\textbar R'\,'\,'(x)\textbar{}\le M\_3. \tag{TB32} {]}

Substitution into the exact formula TB7 gives

{[}

\begin{aligned}
e_f(x)=e^{-8x^2}\bigl\{&
[(256x^2-1089/4)R-32xR'+R'']\cos(16x)\\
&+[512xR-32R']\sin(16x)\bigr\},
\end{aligned}
\tag{TB33}

{]}

and therefore

{[} \boxed{|e_f(x)|\le M_2e^{-8x^2}E(|x|),\quad
E(r)=1+32r+\frac{1345}{8}r^2+256r^3+128r^4.} \tag{TB34} {]}

For example the quadratic coefficient is exactly (1089/8+32=1345/8),
retaining both the original (-1089/4) filter coefficient and the
derivative term (-32xR').

Differentiating TB33 gives the exact coefficients

{[}

\begin{aligned}
e_f'(x)=e^{-8x^2}\bigl\{&
[(13060x-4096x^3)R+(768x^2-3265/4)R'
-48xR''+R''']\cos(16x)\\
&+[(4868-12288x^2)R+1536xR'-48R'']\sin(16x)\bigr\}.
\end{aligned}
\tag{TB35}

{]}

Using all four TB32 bounds yields

{[} \boxed{|e_f'(x)|\le e^{-8x^2}[M_2E_1(|x|)+M_3],} \tag{TB36} {]}

{[}
E\_1(r)=48+\frac{3457}{4}r+3970r\textsuperscript{2+7298r}3+6144r\textsuperscript{4+2048r}5.
{]}

Alternatively (H'\,'\,`(0)=0) and
(\textbar H\^{}\{(4)\}\textbar{}\le M\_4) give
(\textbar R'\,'\,'(x)\textbar{}\le M\_4\textbar x\textbar); the final
(M\_3) in TB36 may then be replaced by (M\_4\textbar x\textbar) when
that yields a smaller integrated bound.

Let (k\_0=f\_0\emph{f\_0) denote the comparison convolution only in this
section, and let (\Delta=k-k\_0). Since
(\Delta=f\_0}e\_f+e\_f\emph{f\_0+e\_f}e\_f), Cauchy--Schwarz gives the
global bounds

{[}

\begin{aligned}
|\Delta(x)|&\le2\|f_0\|_2\|e_f\|_2+\|e_f\|_2^2=:\delta_0,\\
|\Delta''(x)|&\le2\|f_0'\|_2\|e_f'\|_2+\|e_f'\|_2^2=:\delta_2.
\end{aligned}
\tag{TB37}

{]}

All functions are real and even except their odd derivatives, so
(\Delta'(0)=0). Consequently

{[} \textbar{}\Delta(x)-\Delta(0)\textbar{}\le\frac{\delta_2x^2}{2},
\qquad \textbar{}\Delta(x)-\Delta(0)\textbar{}\le2\delta\_0. \tag{TB38}
{]}

These bounds are statements about the exact residual; the original (H),
(f), and (k) remain those of TB1 and TB6. Every norm bound supplied by
TB34 or TB36 is a finite sum of exact Gaussian moments. Specifically,
for integers (m\ge0),

{[} \int\_\{\mathbb R\}\textbar x\textbar\^{}m
e\textsuperscript{\{-16x}2\},dx
=16\^{}\{-(m+1)/2\}\Gamma\left(\frac{m+1}{2}\right). \tag{TB39} {]}

This follows from (v=16x\^{}2) on each half-line; the Gamma factor and
the full power of 16 have both been retained. Squaring the polynomials
in TB34 and TB36 and applying TB39 yields computable certified bounds
for (\textbar e\_f\textbar\_2) and (\textbar e\_f'\textbar\_2). No sign
for the full pairing is claimed by these estimates.

\subsection{7. Certified tails for direct evaluation of the original
moments}\label{certified-tails-for-direct-evaluation-of-the-original-moments}

For each (j\in\{0,2,3,4\}), retain the exact finite integral

{[}
M\_j\textsuperscript{\{{[}16,3{]}\}=\int\_0}3u\textsuperscript{je}\{u\^{}2/32\}
\sum\_\{n=1\}\textsuperscript{\{16\}(2\pi}2n\textsuperscript{4e}\{9u\}-3\pi n\textsuperscript{2e}\{5u\})
e\^{}\{-\pi n\textsuperscript{2e}\{4u\}\},du. \tag{TB40} {]}

We prove, with ample margin,

{[} \boxed{0<M_j-M_j^{[16,3]}<10^{-300}.} \tag{TB41} {]}

All original summands are positive, so the left inequality holds. The
omitted domain is contained in the union of all indices with (u\ge2) and
the indices (n\ge17) with all (u\ge0). Bounding those two sets
separately is sufficient even though their intersection is counted
twice.

First let (u=2+v), (v\ge0). Since (e\^{}3\textgreater20),
(e\textsuperscript{8\textgreater e}6\textgreater400), and hence

{[}
e\textsuperscript{\{4u\}=e}8e\textsuperscript{\{4v\}\textgreater400(1+4v+8v}2).
{]}

Keeping every term in the original exponent before estimating it, we
obtain

{[}

\begin{aligned}
\frac{u^2}{32}+9u-\pi n^2e^{4u}
&\le\frac{145}{8}-1200n^2
-\left(4800n^2-\frac{73}{8}\right)v
-\left(9600n^2-\frac1{32}\right)v^2\\
&\le-1180n^2-v.
\end{aligned}
\tag{TB42}

{]}

The last line uses (145/8\le20n\^{}2), (4800n\^{}2-73/8\ge1), and
positivity of the final quadratic coefficient. For (j\le4),
((2+v)\textsuperscript{j\le(2+v)}4), and

{[} \int\_0\textsuperscript{\infty(2+v)}4e\^{}\{-v\},dv
=16+32+48+48+24=168. {]}

The same majorant TB2 therefore yields

{[} \int\emph{2\^{}\infty u\textsuperscript{je}\{u\^{}2/32\}\Phi(u),du
\le3360\sum}\{n\ge1\}n\textsuperscript{4e}\{-1180n\^{}2\}. {]}

Since (n\^{}2\ge1+3(n-1)), put (q=e\^{}\{-3540\}\textless1/2).
Differentiating the geometric series gives

{[} \sum\_\{n\ge1\}n\textsuperscript{4q}\{n-1\}
=\frac{1+11q+11q^2+q^3}{(1-q)^5}\textless300. {]}

It follows that

{[} \int\_2\^{}\infty u\textsuperscript{je}\{u\^{}2/32\}\Phi(u),du
\textless1008000e\textsuperscript{\{-1180\}\textless10}\{-383\}.
\tag{TB43} {]}

For the last numerical inequality,
(e\textsuperscript{\{1180\}\textgreater e}\{1170\}\textgreater(20)\textsuperscript{\{390\}\textgreater10}\{390\})
and (1008000\textless10\^{}7). This also bounds the smaller integral
from (u=3).

For the index tail, the first two lines of TB3 applied only to (n\ge17)
give

{[}

\begin{aligned}
&\int_0^\infty u^je^{u^2/32}
\sum_{n\ge17}(2\pi^2n^4e^{9u}-3\pi n^2e^{5u})
e^{-\pi n^2e^{4u}}\,du\\
&\hspace{2em}\le\frac{20j!}{3^{j+1}}\sum_{n\ge17}n^2e^{-3n^2}.
\end{aligned}
\tag{TB44}

{]}

Write (n=17+a), (a\ge0) an integer. Then
(n\textsuperscript{2=289+34a+a}2\ge289+35a). With
(q=e\^{}\{-105\}\textless1/2),

{[}

\begin{aligned}
\sum_{n\ge17}n^2e^{-3n^2}
&\le e^{-867}\sum_{a\ge0}(17+a)^2q^a\\
&=e^{-867}\left[\frac{289}{1-q}
+\frac{34q}{(1-q)^2}+\frac{q(1+q)}{(1-q)^3}\right]\\
&<652e^{-867}.
\end{aligned}

{]}

For each requested (j), (20j!/3\^{}\{j+1\}\le20/3), so TB44 is less than
(10000e\textsuperscript{\{-867\}\textless10}\{-369\}). Indeed,

{[}
e\textsuperscript{\{867\}\textgreater20}\{289\}=2\textsuperscript{\{289\}10}\{289\}
\textgreater2\textsuperscript{\{280\}10}\{289\}\textgreater10\textsuperscript{\{84\}10}\{289\}=10\^{}\{373\},
{]}

using (2\textsuperscript{\{10\}\textgreater10}3). Together TB43 and TB44
prove
(M\_j-M\_j\textsuperscript{\{{[}16,3{]}\}\textless10}\{-383\}+10\textsuperscript{\{-369\}\textless10}\{-300\}),
which is TB41. Thus adding a symmetric ball of radius (10\^{}\{-300\})
to a rigorous enclosure of TB40 encloses the exact original moment. No
estimate of the finite integral itself is substituted for its certified
integration.

\subsection{8. Independent derivation of the comparison convolution and
certificate
checks}\label{independent-derivation-of-the-comparison-convolution-and-certificate-checks}

Retain the comparison source from TB31,
(h\_0(x)=Ce\textsuperscript{\{-8x}2\}\cos(16x)), so
(f\_0=(D\^{}2-1/4)h\_0). The cosine product identity and TB14 give

{[}

\begin{aligned}
(h_0*h_0)(x)
&=\frac{C^2}{2}e^{-4x^2}
\int_{\mathbb R}e^{-16z^2}[\cos(32z)+\cos(16x)]\,dz\\
&=\frac{C^2\sqrt\pi}{8}e^{-4x^2}
[e^{-16}+\cos(16x)].
\end{aligned}
\tag{TB45}

{]}

The Gaussian Fourier integral here is
(\int e\textsuperscript{\{-16z}2\}\cos(32z),dz=(\sqrt\pi/4)e\^{}\{-16\}),
with its full exponential factor retained. Since all sources and their
derivatives have Gaussian decay, differentiating under convolution and
integrating by parts give

{[} k\_0=f\_0\emph{f\_0=(D\textsuperscript{2-1/4)}2(h\_0}h\_0). {]}

For an algebraic verification, put (q=-8x+16i). For
(a(x)=e\textsuperscript{\{-4x}2+16ix\}),

{[} D\textsuperscript{2a=(q}2-8)a,\qquad
D\textsuperscript{4a=(q}4-48q\^{}2+192)a, {]}

and hence

{[} (D\textsuperscript{2-1/4)}2a
=\left(q\textsuperscript{4-\frac{97}{2}q}2+\frac{3137}{16}\right)a. {]}

The same identity with (q=-8x) applies to (e\textsuperscript{\{-4x}2\}).
Expanding the real and imaginary parts proves exactly

{[} \boxed{\begin{aligned}
k_0(x)=\frac{C^2\sqrt\pi}{128}e^{-4x^2}\bigl[&
(65536x^4-1622528x^2+1250369)\cos(16x)\\
&+2048x(256x^2-1121)\sin(16x)\\
&+e^{-16}(65536x^4-49664x^2+3137)\bigr].
\end{aligned}} \tag{TB46} {]}

This is the exact expression used in \texttt{certify\_pairing.py};
\texttt{K0\_EXACT.txt} contains the same expression with its common
factor (C\^{}2) handled by the script.

The polynomial envelopes for the comparison function and derivative are

{[}

\begin{aligned}
|f_0(x)|&\le Ce^{-8x^2}A(|x|),
&A(r)&=\frac{1089}{4}+512r+256r^2,\\
|f_0'(x)|&\le Ce^{-8x^2}B(|x|),
&B(r)&=4868+13060r+12288r^2+4096r^3.
\end{aligned}
\tag{TB47}

{]}

TB31 proves the first bound; TB35 with (R=C) constant proves the second
derivative formula, or it follows by direct differentiation of TB31.
TB39 computes the squared envelope norms exactly. Since (C=M\_0\le1/3),
the envelope for (f\_0) is bounded by the envelope TB8 used for (f).
Therefore the universal convolution and Gamma-tail bounds TB16 and TB29
apply to (k\_0) as well.

For clarity, the script's additional error expressions have the
following rigorous justifications.

For (0\textless{}\varepsilon\le1), the omitted origin interval of the
Gamma integral of (k\_0) has absolute value at most

{[} \frac32k\_0(0)\varepsilon +\frac{\|f_0'\|_2^2\varepsilon^2}{8}.
\tag{TB48} {]}

This is the integration of the TB25 majorant using
(0\le k\_0(0)-k\_0(x/2)\le\textbar f\_0'\textbar\_2\textsuperscript{2x}2/8)
and (1-e\^{}\{-x\}\ge x/2). Any upper bound for
(\textbar f\_0'\textbar\_2), such as TB47, is valid in TB48.

Let (\Delta=k-k\_0), with the exact residual bounds
(\delta\_0,\delta\_2) from TB37. For (0\textless h\le1), the complete
Gamma functional difference satisfies

{[}

\begin{aligned}
|A_\infty(k)-A_\infty(k_0)|
\le{}&\left[\gamma_{\rm E}+\log\pi+\frac32
+4\log\frac1h
+\frac{e^{-1}+4e^{-1/4}}{1-e^{-1}}\right]\delta_0
+\frac{\delta_2h^2}{8}.
\end{aligned}
\tag{TB49}

{]}

To verify every interval in this bound, on ((0,1{]}) write the Gamma
numerator difference exactly as

{[} (e\textsuperscript{\{-x\}-e}\{-x/4\})\Delta(0)
+e\^{}\{-x/4\}{[}\Delta(0)-\Delta(x/2){]}. {]}

The first part integrates to at most (3\delta\_0/2), by
(\textbar e\textsuperscript{\{-x\}-e}\{-x/4\}\textbar{}\le3x/4) and
(1-e\^{}\{-x\}\ge x/2). For the second part use
(\textbar{}\Delta(0)-\Delta(x/2)\textbar{}\le\delta\_2x\^{}2/8) on
((0,h{]}), yielding (\delta\_2h\^{}2/8), and use (2\delta\_0) on
({[}h,1{]}), yielding (4\delta\_0\log(1/h)). On ({[}1,\infty)) use the
original numerator directly to get
(\delta\_0\int\emph{1\textsuperscript{\infty(e}\{-x\}+e\textsuperscript{\{-x/4\})/(1-e}\{-1\}),dx).
Its value is the last coefficient in TB49. The retained prefactor
contributes ((\gamma}\{\rm E\}+\log\pi)\delta\_0). This proves the
script's full Gamma error at (h=1/16).

Finally, for (x\ge0), define the exact envelope integral

{[}
\mathcal C\_\{P,Q\}(x)=\int\_\{\mathbb R\}e\textsuperscript{\{-16z}2\}
P(\textbar z\textbar+x/2)Q(\textbar z\textbar+x/2),dz, {]}

where (P,Q) have nonnegative coefficients. Under (y=x/2+z), both
(\textbar y\textbar) and (\textbar x-y\textbar) are at most
(\textbar z\textbar+x/2). Therefore TB34 and TB47 give

{[} \textbar{}\Delta(x)\textbar{}\le e\textsuperscript{\{-4x}2\}
{[}2CM\_2\mathcal C\_\{A,E\}(x)+M\_2\^{}2\mathcal C\_\{E,E\}(x){]}.
\tag{TB50} {]}

The script expands these positive polynomials and integrates every power
with TB39. Multiplying TB50 at (x=\log n) by (2\Lambda(n)/\sqrt n) and
summing over all prime powers (2\le n\le64) gives the stated finite
prime error. The remaining original prime tail is bounded by TB23.

Thus the exact-convolution formula, original-source moment tails,
derivative residuals, finite prime error, origin error, and Gamma error
used by the certificate all have independent proofs here. This audit
inspected the mathematics and script expressions; it did not rerun the
numerical Weil-pairing calculation or independently assert its final
numerical enclosure.

\subsection{Derivation log}\label{derivation-log}

\begin{enumerate}
\def\labelenumi{\arabic{enumi}.}
\tightlist
\item
  Read the programme source file, with particular attention to the exact
  variables, heat time, filter, convolution, prime factor two, and Gamma
  numerator.
\item
  Derived the all-orders moment bound TB5 from the complete original
  density by positive majorants; kept the source's negative summand
  explicit in TB1--2.
\item
  Expanded the actual differential filter and derived the polynomial and
  Gaussian bounds TB7--17.
\item
  Proved the arithmetic sum-to-integral comparison, retaining all prime
  powers, and the rational cutoff TB23.
\item
  Split the Gamma tail only for a strictly positive lower endpoint,
  evaluated its (k(0)) contribution exactly, and proved TB26--30.
\item
  Independently verified the receiving calculation's residual polynomial
  and supplied the derivative residual, TB31--39.
\item
  Evaluated the displayed rational coefficient identities with Python's
  exact \texttt{fractions.Fraction}; the proofs also specify the finite
  formulas from which each coefficient can be reproduced.
\item
  Proved TB40--44 for the requested original moment truncation, with
  tails strictly smaller than the script's radius (10\^{}\{-300\}).
\item
  Read \texttt{certify\_pairing.py}, \texttt{K0\_EXACT.txt}, and the
  current \texttt{CERTIFICATE.json}; independently derived TB45--50 and
  checked each mathematical error expression used by the script. No
  numerical pairing was rerun.
\end{enumerate}

No publication, remote mutation, rendering, source indexing, numerical
Weil-pairing evaluation, or assertion of its sign was performed by this
derivation.

\end{document}
